% Este arquivo .tex será incluído no arquivo .tex principal. Não é preciso
% declarar nenhum cabeçalho

\section{Agora eu tenho um e-mail da Unicamp}

Ao ingressar num curso de computação, você recebe pelo menos duas contas de
e-mail: do IC e da DAC. Estes são os principais meios de comunicação da
Universidade com você, portanto fique esperto e não deixe de ler esses emails
regularmente!

Para acessar o webmail do IC, o endereço é \url{webmail.students.ic.unicamp.br}.
O usuário e a senha são os mesmos do sistema Linux do IC, que você receberá nas
primeiras semanas de aula. Caso tenha dúvidas, dê uma passada na Secretaria de
Graduação, que fica no IC-2 (prédio ao lado das Artes Cênicas e para cima da
Economia) e pergunte!

Uma dica interessante, que muitas vezes passa despercebida, é que no mesmo
documento em que você recebe sua senha, vem indicado um {\it alias} para o seu
e-mail. Assim, você poderá utlizá-lo de uma forma mais amigável, ao invés de ser
somente o seu próprio RA, por exemplo:

\begin{center}
\texttt{francisco.silva@students.ic.unicamp.br}\\
em vez de\\
\texttt{ra129873@students.ic.unicamp.br}
\end{center}

A DAC também dispõe um webmail para o aluno. O login é a primeira letra do nome
do aluno, seguido dos dígitos do RA. O e-mail da DAC também é útil. É nele que
os alunos são avisados sobre o período de matrícula, recebem os seus pedidos de
matrículas provisórias, avisos de desistências e trancamentos, além de informar
os alunos a respeito de eventos que acontecem na Unicamp, como feiras,
palestras, festivais, eleições. O endereço do webmail da DAC é
\url{webmail.dac.unicamp.br}.

Para a galera da Engenharia, ainda há o da FEEC, que muitos sequer ficam
sabendo que existe! Só um semestre ou até um ano depois vão até o SIFEEC retirar
seu login e senha. Ou então ficam sabendo só no segundo ou terceiro ano de curso
e aí há mais de 700 e-mails não lidos. Assim como o do IC, é muito utilizado
para divulgação de eventos, oportunidades de estágios e de iniciação científica.
Para utilizá-lo, você deve ir até a FEEC e procurar o SIFEEC, que é o local
responsável por isso. Fica no segundo andar do prédio de laboratórios da FEEC
(um com escadas amarelas). O endereço do webmail da FEEC é
\url{webmail.fee.unicamp.br} e a senha é a mesma do Linux da FEEC.

\begin{figure}[b!]
    \centering
    \includegraphics[width=.3\textwidth]{img/alem_da_graduacao/email.jpg}
\end{figure}
Você também pode redirecionar os e-mails que receber nas contas do IC, FEEC e
DAC para qualquer outra conta (no Gmail por exemplo).

\subsection{Redirecionamento de e-mails}
Para efetuar o redirecionamento do e-mail institucional para outro e-mail, siga
os passos abaixo para cada e-mail institucional que você tiver.

\subsubsection*{E-mail da DAC}

\begin{compactenumerate}
    \item  Acesse \url{www.dac.unicamp.br}
    \item  Acesse Serviços On-line
    \item  Acesse Alunos
    \item  Clique em Acesso aos Serviços Acadêmicos
    \item  Faça login
    \item  Clique em Redirecionamento de Email
\end{compactenumerate}

\subsubsection*{E-mail da FEEC}

\begin{compactenumerate}
    \item  Acesse \url{webmail.fee.unicamp.br}
    \item  Faça login
    \item  Acesse Options
    \item  Acesse Mail Forwarding
\end{compactenumerate}

\subsubsection*{E-mail do IC}

\begin{compactenumerate}
    \item  Acesse \url{webmail.students.ic.unicamp.br}
    \item  Faça login
    \item  Acesse Filtros
    \item  Clique em Adicionar uma nova regra
    \item  Em Condição, na primeira lista drop-down com a opção Header, escolha a opção All
    \item  Em Ação, escolha Redirecionar para o seguinte endereço de e-mail
\end{compactenumerate}

\subsection{Suporte técnico}

Além dos serviços de e-mails, ao entrar na Unicamp, você também ganha acesso a
diversos outros serviços, como a VPN da Unicamp, acesso remoto ao IC etc.

Obviamente configurar essas coisas não é trivial, o IC mantem uma página com
instruções de como usar os seus serviços comuns, como configurar seu cliente de
email para acessar o servidor do IC e acessar o IC remotamente.

Sempre que precisar de ajuda no IC, confira se não há uma resposta para seu
problema nesta página: \url{suporte.ic.unicamp.br/index.php/Alunos}.

% TODO: Achar coisas sobre o suporte da FEEC e da DAC

\subsection{Tenha uma conta no Gmail!}

Nós sabemos que é difícil. Muitas vezes você vem usando um serviço de e-mail
durante anos, seja ele Hotmail, Yahoo, Uol, Terra, Zipmail, e mudar é
trabalhoso.  Mas confie na gente: agora vai ser mais fácil. Você está começando
uma vida nova na universidade, quando já tiver cadastrado seu endereço antigo em
vários serviços acadêmicos e divulgado entre os novos contatos que vai fazer,
será bem mais complicado.

Entre as vantagens do Gmail, podemos destacar:

\begin{itemize}
    \item Muito espaço, nunca mais apague nada.
    \item Visualização em threads, útil para acompanhar discussões.
    \item Bom aplicativo para smartphones.
    \item Filtros e marcadores infinitos para te ajudar.
    \item Uma pesquisa que funciona.
    \item Interface mais polida que as da concorrência.
\end{itemize}

Qualquer dúvida, novamente, procure um veterano!

Finalizando, não deixe de estar sempre informado sobre os acontecimentos ou
divulgações da Unicamp, do CACo, da AAACEC e da Conpec, essas três entidades
compostas por alunos de computação.
